% =============================================================================
%  正文 1：§1 引言、§2 GL(n,Q) 的猜想、§3 域
% =============================================================================

\section{引言}
\label{sec:intro}

本文是 2004 年 3 月 8 日我在高等研究院（I.A.S.）给一群物理学家所作几次讲座的扩写版。

这是一份\textbf{进行中的工作}：过几个月再回来看，末尾那些空白的小节也许就补齐了。
本文分两层写。原始讲座里没有讲、第一遍阅读可以略过的许多技术细节，都收在尾注里。
每一节的头几段，都写得让尽可能多的读者能读下去。

本文充其量只是许多优秀综述——[Ar, F, G1, G2, Gr, Kn, K, R1, T]——的导引，
这些综述讲自守形式与朗兰兹纲领。
再深一点的综述有 [BR]，专著有 [Ba, Be]，评论性文章有 [M, R2, R3]。

\begin{yizhu}[为什么「为物理学家而写」值得单列]
作者面对的听众是一群在 I.A.S. 工作的物理学家：他们熟悉规范场论、路径积分与 $\zeta$ 函数，
但没有受过数论的训练。全文的写法因此有一个特点：\textbf{凡是能用一个具体公式说清的地方，就不写一般性定义}。
例如第 5 节讲 Hecke 算子，直接画对应关系、直接给积分公式，而不去定义 Hecke 范畴。
物理读者顺着第 4 节的 $L$ 函数读下去，往往比数论读者还快。
\end{yizhu}

\section{$\GL(n,\bbQ)$ 的猜想}
\label{sec:gl-n}

粗略地说，猜想是说：应该存在一个对应
\begin{equation*}
\begin{Bmatrix}
\text{较好的 $n$ 维不可约表示}\\
\Gal(\Qbar/\bbQ)\ \text{的}
\end{Bmatrix}
\;\longrightarrow\;
\begin{Bmatrix}
\text{较好的自守表示}\\
\GL(n,\bbA_{\bbQ})\ \text{的}
\end{Bmatrix}
\tag{2.0.1}
\end{equation*}
使得
\begin{equation*}
\{\text{Frobenius 的特征值}\}\;\longrightarrow\;\{\text{Hecke 算子的特征值}\}\ .
\tag{2.0.2}
\end{equation*}

接下来几节的目的，是把这句话里的每个词讲清楚。
讲完之后，我们再简单地看一下这个论断向其它域（$\bbQ$ 以外）与其它代数群（$\GL(n)$ 以外）的种种推广。

\begin{yizhu}[「较好的」是什么意思]
原文写作 nice，作者通篇没有定义它。这不是疏忽：\textbf{定义「较好」正是朗兰兹纲领的全部内容}。
粗略地说，左边要求表示在某种意义下「来自几何」——例如来自某个簇的平展上同调，
从而带有 Frobenius 特征值、满足 Weil 猜想所要求的那类整除性；
右边要求表示是「尖点」的、并且满足某种解析增长条件。
两边各自的好条件必须恰好对上，对应才可能是一一的（第 7 节会看到 $n=2$ 时确实对上了）。
\end{yizhu}

\section{域}
\label{sec:fields}

\subsection{代数簇上的点}

设 $E\subset F$ 是域。Galois 群 $\Gal(F/E)$ 是 $E$ 的域自同构 $\varphi:E\to E$ 中固定 $F$ 每个元素的那些。

\begin{yizhu}[原文疑误：$\Gal(F/E)$ 的方向]
原文此处写作「$\Gal(F/E)$ 是 $E$ 的域自同构 $\varphi:E\to E$ 中固定 $F$ 每个元素者」，
方向恰好颠倒。按标准定义，$\Gal(F/E)$ 应是 $F$ 的域自同构（$\varphi:F\to F$）中固定 $E$ 每个元素者。
译文照录原文以便对读；后文（例如 $\Gal(\overline{\bbQ}_p/\bbQ_p)$、$\Gal(\overline{\bbF}_q/\bbF_q)$）
用的都是标准方向，读者按标准理解即可。
\end{yizhu}

设 $E[x_1,x_2,\dots,x_n]$ 是以 $E$ 为系数的 $n$ 元多项式代数。
若 $f_1,f_2,\dots,f_r\in E[x_1,x_2,\dots,x_n]$ 是多项式，记 $X(E)$ 为解集
\[
\{x\in E^n\mid f_1(x)=f_2(x)=\cdots=f_r(x)=0\}.
\]
我们也可以把 $f_i$ 看成系数在 $F$ 中的多项式，于是坐标在 $F$ 中的解集 $X(F)$ 也有意义。
此时 $\Gal(F/E)$ 作用在 $X(F)$ 上，并固定 $X(E)$。

这时称 $X(F)$ 为代数簇 $X$ 在域 $F$ 中的点集。
我们说代数簇 $X$ 同时表示了它在 $E$ 的所有扩域上的解集，但说 $X$ 是\textbf{定义在 $E$ 上的}。
当 $f_i$ 是齐次多项式、$X(E)$ 是射影空间 $\bbP^{r-1}(E)$ 中相应的解集时，同样的说法照搬。

特别地，若代数簇 $X$ 定义在整数环 $\bbZ$ 上，那么它同时活在两个世界里：
$X(\bbC)$ 是一个复代数簇，而 $X\pmod p$ 是定义在 $\bbF_p$ 上的代数簇（把 $X$ 的方程模 $p$ 约化即得）。
不仅如此，对任何扩域 $\bbF_{p^r}$，都可以考虑坐标在 $\bbF_{p^r}$ 中的点集 $X\pmod p(\bbF_{p^r})$。
这两个世界之间惊人的联系，将在 \S\ref{sec:vary-p} 中说明。

\subsection{有限域}

对每个素数 $p$ 与每个正整数 $n$，存在唯一一个含 $p^n$ 个元素的域 $\bbF_{p^n}$。
若 $n=1$，则 $\bbF_p=\bbZ/(p)$ 就是模 $p$ 的整数。
若 $p_1,p_2$ 是素数，则 $\bbF_{p_1^n}\subset\bbF_{p_2^m}$ 当且仅当 $p_1=p_2$ 且 $n$ 整除 $m$。
现在设 $q=p^n$（$p$ 为素数），则对每个 $x\in\bbF_q$ 有
\[
px=0,\qquad x^q=x.
\]
我们说 $\bbF_q$ 的特征为 $p$。
Galois 群 $\Gal(\bbF_{q^r}/\bbF_q)$ 同构于 $\bbZ/(r)$，由 Frobenius 映射 $\sigma_q(x)=x^q$ 生成。
它是自同构，因为
\[
(x+y)^q=x^q+qx^{q-1}y+\binom{q}{2}x^{q-2}y^2+\cdots+y^q .
\]
其中每个二项式系数都能被 $p$ 整除，故 $(x+y)^q=x^q+y^q$。
当 $r$ 变化时，这些 Frobenius 自同构彼此相容。取 $p=q$ 就得到 $\Gal(\overline{\bbF}_p/\bbF_p)$ 中一个特定的元素 $\sigma_p$。

\subsection{数域}

数域是 $\bbQ$ 的有限扩张。每个数域都可以由在 $\bbQ$ 上添加某个有理系数多项式的根得到。
代数闭包 $\overline{\bbQ}\subset\bbC$ 是所有数域的「并」：它是所有有理系数多项式的所有根组成的集合。
正如有理数 $\bbQ$ 由所有分数 $a/b$（$a,b\in\bbZ$）组成，任何数域 $E$ 都由所有分数 $a/b$ 组成，
其中 $a,b\in\sO_E$，而 $\sO_E$ 是 $E$ 中的整数环。
例如 $E=\bbQ[i]$ 时，$\sO_E$ 就是所有 $a+bi$（$a,b\in\bbZ$）。

\subsection{局部函数域}

设 $\bbF_q$ 是有限域，$\bbF_q[[t]]$ 是以 $\bbF_q$ 为系数的形式幂级数环，加法与乘法按显然的方式定义。
元素 $a=\sum_{n\ge 0}a_nt^n$ 可逆当且仅当 $a_0\neq 0$。
此时逆元可以这样求：设 $b=\sum_{m\ge 0}b_mt^m$，把 $ab=1$ 展开，
再对系数 $b_m$ 解出线性方程组。
映射 $\bbF_q[[t]]\to\bbF_q$ 把上述元素 $a$ 送到常数项 $a_0$，它是满的环同态，
因此称 $\bbF_q$ 为\textbf{剩余域}。

域 $\bbF_q((t))$ 由所有分数 $a/b$（$a,b\in\bbF_q[[t]]$）组成，我们称 $\bbF_q[[t]]$ 为域 $\bbF_q((t))$ 中的整数环。
容易看出，任何 $A\in\bbF_q((t))$ 都可以写成形式 Laurent 级数
\[
A=\sum_{n=N}^{\infty}A_nt^n
\]
（$N\in\bbZ$ 可为负），其中只有有限多项含 $t$ 的负幂。
这样的元素 $A$ 的\textbf{赋值}，是使 $A_n\neq 0$ 的最小整数 $n$。

类似地可以构造函数域 $\bbC((t))$ 及其整数环 $\bbC[[t]]$；
不过这种情形对数论学家意义不大（对物理学家意义更大）。

\subsection{$p$ 进域}

固定素数 $p$。一个 $p$ 进整数是形式幂级数 $a=a_0+a_1p+a_2p^2+\cdots$，其中 $0\le a_i\le p-1$。
这样的元素全体记作 $\bbZ_p$。加法与乘法用幂级数运算，但要「进位」：
每当遇到大于 $p$ 的系数 $a$，就把它写成 $a_0+a_1p$，留下 $a_0$ 部分，把 $a_1$ 部分进到下一位。
因此，映射 $a\mapsto a_0$ 定义了一个满的环同态 $\bbZ_p\to\bbZ/(p)$，即到剩余域的同态。
环 $\bbZ_p$ 包含了大部分有理数：它显然包含整数，但只要求 $p\nmid b$，分数 $a/b$ 也在其中。
例如在 $\bbZ_5$ 中，$3$ 的逆是
\[
3^{-1}=2+3\cdot 5+1\cdot 5^2+3\cdot 5^3+1\cdot 5^4+\cdots,
\]
把右边乘以 $3$ 即可验证。

然而 $5$ 在 $\bbZ_5$ 中不可逆，所以要得到域，必须把它的逆、$5^2$ 的逆等等都添进来。
这给出 $\bbQ_p$ 的两种描述：一是所有分数 $a/b$（$a,b\in\bbZ_p$），二是所有形式 Laurent 级数
\[
c=\sum_{n=N}^{\infty}c_np^n
\]
（只有有限多个含 $p$ 的负幂）。元素 $c$ 的赋值 $\val(c)$ 是展开式中首项的「次数」。
特别地，$\bbQ_p$ 包含 $\bbQ$。

$p$ 进数也可以实现为有理数 $\bbQ$ 关于某个范数 $|\ |_p$ 的（拓扑）完备化\en{1}。
有了这个度量空间结构，域 $\bbQ_p$ 是局部紧的。

注意 $\bbQ_p$ 的特征为 $0$，而且可以把 $\bbQ_p$ 嵌入复数域 $\bbC$——从现在起，我们假定已经这样做了\en{3}。

$\bbQ_p$ 有许多有限扩张，都称为 $p$ 进域。每个这样的域都可以通过把数域 $E$ 在某个包含（或「位于」）$p$ 的素理想 $\mathfrak p$ 处「完备化」得到。
$p$ 进域 $E_{\mathfrak p}$ 含有一个整数环 $\sO$，它再投影到剩余域 $\sO_E\to\bbF$，这是 $\bbF_p$ 的有限扩张。
若 $\deg(E_{\mathfrak p}/\bbQ_p)=\deg(\bbF/\bbF_p)$，则称 $E_{\mathfrak p}$ 是 $\bbQ_p$ 的\textbf{非分歧扩张}。
$\bbQ_p$ 在每个次数上都有唯一的非分歧扩张，但也还有别的扩张。
$p$ 进域统称\textbf{局部域}。

\begin{yizhu}[两种「局部」与两种「整体」]
这里的「局部/整体」对应，与几何里的局部化是同一件事：整体域（$\bbQ$、数域）像一个函数域，
局部域（$\bbQ_p$、$\bbF_q((t))$）像它在一个点处的形式邻域。
第 6 节的平代环 $\bbA_{\bbQ}=\bbR\times\prod'_p\bbQ_p$，就是把所有「点」处的信息加一个「无穷远点」$\bbR$ 打包在一起。
几何朗兰兹把这套语言整体几何化：数域换成曲线 $X$，局部域换成 $X$ 的完形局部域，平代环换成 $X$ 上所有点的乘积。
\end{yizhu}

\subsection{整体函数域}

设 $Y$ 是紧 Riemann 面，则 $Y$ 上的亚纯函数 $\bbC(Y)$ 构成一个域。
这样的 $Y$ 具有（光滑）复射影代数簇的结构，可以实现为 $\bbP^2(\bbC)$ 的子簇。
因此它有「有理函数」域，而这也正是 $\bbC(Y)$。
若 $Y=\bbP^1(\bbC)$，其函数域就是「有理函数」域 $\bbC(x)$，由商 $p(x)/q(x)$（$p,q$ 为多项式）组成。
若 $Y'\to Y$ 是紧 Riemann 面的（可能分歧的）有限覆叠，则 $\bbC(Y')$ 是 $\bbC(Y)$ 的有限扩张。

类似地，设 $Y$ 是定义在有限域 $\bbF(q)$ 上的光滑射影代数曲线，$\bbF_q(Y)$ 是其（「有理」）函数域。
若 $Y=\bbP^1$，则 $\bbF_q(Y)=\bbF_q(x)$ 就是商 $p(x)/q(x)$ 的集合，其中 $p,q\in\bbF_q[x]$。
对任何一点 $y\in Y(\bbF_q)$，在 $y$ 处选定坐标 $t_y$ 就决定了 $\bbF_q(Y)$ 的一个完备化，
它同构于局部函数域 $\bbF_q((t))$。正因如此，$y\in Y$ 这些点被称为\textbf{位}（place）。

\subsection{高维的域}

上述这些域（连同它们的有限扩张与各种代数闭包），基本上穷尽了人们传统上关心的域。
但可以用两种构造从旧域 $k$ 不断造出新域 $K$：
（a）把整数环 $\sO_k$ 关于某个范数完备化；
（b）取函数域，$K=k((t))$ 或 $K=k(Y)$。

于是，例如有函数域
\[
\bbQ_p(Y)((t)),
\]
它由以 $\bbQ_p(Y)$ 为系数的形式 Laurent 多项式组成，而 $\bbQ_p(Y)$ 是定义在 $\bbQ_p$ 上的 $n$ 维代数簇 $Y$ 上的 $p$ 进函数域。

\subsection{全局图景}

总而言之，数论学家最关心的那些域可以排进一张表
（其中 $E$ 是数域，$E_{\mathfrak p}$ 是把 $E$ 在素理想处完备化得到的 $p$ 进域；$p$ 是素数，$q=p^r$ 是 $p$ 的幂）。

\begin{center}
\renewcommand{\arraystretch}{1.35}
\begin{tabular}{@{}lcccc@{}}
\toprule
& \multicolumn{2}{c}{基域} & \multicolumn{2}{c}{扩域} \\
\cmidrule(lr){2-3}\cmidrule(lr){4-5}
整体 & $\bbQ$ & $\bbF_p(X)$ & $E$ & $\bbF_q(X)$ \\
局部 & $\bbQ_p$ & $\bbF_p((t))$ & $E_{\mathfrak p}$ & $\bbF_q((t))$ \\
整数 & $\bbZ_p$ & $\bbF_p[[t]]$ & $\sO_{\mathfrak p}$ & $\bbF_q[[t]]$ \\
剩余 & $\bbF_p$ & $\bbF_p$ & $\bbF_q$ & $\bbF_q$ \\
\bottomrule
\end{tabular}
\end{center}

数域 $\bbQ$ 与 $E$ 在「无穷远」处也有完备化，给出局部域 $\bbR$（作为 $\bbQ$ 唯一的无限完备化）
以及 $\bbR,\bbC$（作为 $E$ 可能的无限完备化）。

\begin{keytakeaways}
这一节的用途是建立一张对照表：\textbf{凡是数论里出现的「域」，在几何里都有对应的「曲线上的点」}。
整体域 $\leftrightarrow$ 曲线；局部域 $\leftrightarrow$ 曲线上的点；平代环 $\leftrightarrow$ 曲线所有点的乘积。
第 6 节把自守表示写在平代环上；几何朗兰兹则把它写在曲线的 Hecke 栈上。
记住这张表，后面所有「推广到其它域」的句子，都可以读成「把曲线换掉」。
\end{keytakeaways}
